%webamc CODE=r401/2023/vpn2
%webamc TITLE=Tunnel VPN de niveau 2
%webamc RND=0

Soit le réseau ci-dessous avec les IP associées aux interfaces:

\begin{center}
  \begin{minipage}{0.69\textwidth}\begin{tikzpicture}
  \machine{}{C1}{270}{client1}
  \node[right of=C1,node distance=1cm] (pt) {};
  \serveur{right of=C1,node distance=2cm}{V1}{270}{vpn1}
  \petitlan{above of=pt,node distance=1.5cm}{L1}{\scriptsize LAN1}
  \routeur{right of=L1,node distance=2cm}{R1}{90}{routeur1}
  \draw[lien] (R1) -- node[pos=0.2,above] {\scriptsize eth0} (L1);
  \draw[lien] (C1) -- node[pos=0.2,left] {\scriptsize eth0} (L1);
  \draw[lien] (V1) -- node[pos=0.2,right] {\scriptsize eth0} (L1);
  %
  \lan{right of=R1,node distance=2cm}{I}{\scriptsize 1.2.3.0/24}
  \routeur{right of=I,node distance=2cm}{R2}{90}{routeur2}
  \draw[lien] (R1) -- node[pos=0.2,above] {\scriptsize eth1} (I);
  \draw[lien] (R2) -- node[pos=0.2,above] {\scriptsize eth1} (I);
  %
  \petitlan{right of=R2,node distance=2cm}{L2}{\scriptsize LAN2}
  \serveur{right of=V1,node distance=6cm}{V2}{270}{vpn2}
  \machine{right of=V2,node distance=2cm}{C2}{270}{client2}
  \draw[lien] (R2) -- node[pos=0.2,above] {\scriptsize eth0} (L2);
  \draw[lien] (C2) -- node[pos=0.2,right] {\scriptsize eth0} (L2);
  \draw[lien] (V2) -- node[pos=0.2,left] {\scriptsize eth0} (L2);
  \draw[-,ultra thick,red] (V1) -- node[above] {\scriptsize tunnel VPN} (V2);
\end{tikzpicture}
\end{minipage}
  \begin{minipage}{0.3\textwidth}
    \begin{tabular}{|l|l|l|}
      \hline
      Hôte & Interface & IP
      \\
      \hline
      client1 & eth0 & 10.1.0.1\\
      vpn1 & eth0 & 10.1.0.100\\
      routeur1 & eth0 & 10.1.0.254\\
      \hline
      client2 & eth0 & 10.2.0.1\\
      vpn2 & eth0 & 10.2.0.100\\
      routeur2 & eth0 & 10.2.0.254\\
      \hline
    \end{tabular}
  \end{minipage}
\end{center}

On a mis en place un tunnel VPN entre vpn1 et vpn2 pour interconnecter
les clients des deux réseaux locaux LAN1 et LAN2.

On suppose dans cet exercice:
\begin{itemize}
\item que le tunnel est de {\em niveau 2} ;
\item que vpn1 et vpn2 utilisent UDP pour transporter leurs données~;
\item et enfin que le chiffrement n'altère pas la taille des données.
\end{itemize}

Rappels: en-tête + FCS ethernet = 26 octets ; en-tête IP = en-tête TCP
= 20 octets ; et en-tête UDP = en-tête ICMP = 8 octets.
